\documentclass[10pt,a4paper]{amsart}
\usepackage[margin=19mm]{geometry}
\usepackage{amsmath,amssymb,mathtools}
\usepackage[scr=rsfs,cal=euler]{mathalfa}
\usepackage{xfrac}
\usepackage[unicode,hidelinks,bookmarks=false]{hyperref}
\newcommand{\ie}{i.e.\ }
\newcommand{\cf}{cf.\ }
\newcommand{\resp}{respectively\ }
\newcommand{\Subsection}{\subsection}
\providecommand{\nobreakdash}{\nobreak\hbox{-}}
\setlength{\emergencystretch}{2em}
\multlinegap0pt
\usepackage{bm}
\usepackage{bbm}
\usepackage{dsfont}
\usepackage{xspace}
\usepackage{cancel}
\usepackage{mathtools}
\usepackage{amsthm}
\makeatletter
\@ifundefined{thmo}{\newtheorem{thmo}{Thmo}}{}
\@ifundefined{dfn}{\newtheorem{dfn}{Definition}}{}
\@ifundefined{lem}{\newtheorem{lem}[dfn]{Lemma}}{}
\@ifundefined{prop}{\newtheorem{prop}[thmo]{Proposition}}{}
\makeatother
\newcommand{\coveropt}[2]{{\mathcal{C}\mathcal{C}_{#1}(#2)}}
\newcommand{\dC}{{\mathcal{C}}}
\newcommand{\dF}{{\mathcal{F}}}
\newcommand{\dS}{{\mathcal{S}}}
\title[On the $4$-clique cover number of graphs]{On the $4$-clique cover number of graphs}
\author{Yihan Chen, Jialin He, Tianying Xie}
\date{Source: arXiv:2506.10478v1 (CC BY 4.0)}
\begin{document}
\maketitle

\noindent\emph{Source and licence.} On the $4$-clique cover number of graphs. Version arXiv:2506.10478v1 (CC BY 4.0), distributed under the Creative Commons Attribution 4.0 International licence, as recorded in the arXiv licence metadata for this exact version and shown on the abstract page https://arxiv.org/abs/2506.10478v1. Full rights notice: licenses/arxiv-2506.10478-CC-BY-4.0.txt.\par\smallskip

\noindent\textit{Editorial scope.} Counted unit: the Lem whose source label is \texttt{Lem:Bound of 3-clique cover by f(A)} in the article (source0.tex lines 582--587, source label \texttt{Lem:Bound of 3-clique cover by f(A)}), delivered together with the article's own complete original proof of it (source0.tex lines 588--650, one \texttt{proof} environment of 63 lines). No mathematics is written, repaired or re-derived: statement and proof are the article's own text line for line. Required context retained uncounted: Claim:p<n-delta (lines 336--342); Dfn:greedy sequence (lines 564--575). Every definition, display and label of those lines is retained unchanged. Omitted from this package and not redistributed: 1--335, 343--563, 576--581, 651--1368. Nothing inside a retained excerpt is omitted or altered: every retained body is delivered exactly as the article's lines are written, so no omission receipt of any kind accompanies this package. Citations: the retained excerpts carry no citation command at all, so no bibliography entry of the article is transcribed and no reference is redistributed. Reference adaptation: none was needed and none was made; every reference inside the delivered unit resolves inside it, so no unresolved reference or citation is produced. Printed slips kept literal: no printed word, symbol, hypothesis or number of the counted statement or of its proof was altered. Layout and class: the article's own class is \texttt{article} with packages including pgf, tikz, footmisc, verbatim, caption, epsf, epsfig, amsfonts, amsgen, amsmath, amstext, amsbsy, amsopn, amsthm; the wrapper is the lane's pinned \texttt{amsart} editorial wrapper at 10pt on A4, which restates the theorem-like environments the retained text needs and adds the article's own macro definitions that the retained text needs (\texttt{\textbackslash coveropt}, \texttt{\textbackslash dC}, \texttt{\textbackslash dF}, \texttt{\textbackslash dS}), with extra wrapper packages (bm, bbm, dsfont, xspace, cancel, mathtools). The delivered numbering is the unit's own. Assets: no retained span contains a figure environment, a graphics-inclusion command, a PDF-page inclusion, a table or a TikZ picture; no image file is copied into this package, the package declares no asset dependency and no image file is redistributed with it. Licence: version arXiv:2506.10478v1 of this article is distributed under the Creative Commons Attribution 4.0 International licence, as recorded in the arXiv licence metadata for this exact version and shown on the abstract page https://arxiv.org/abs/2506.10478v1. A full rights notice is delivered at licenses/arxiv-2506.10478-CC-BY-4.0.txt. Characters: every retained excerpt is pure ASCII, so the delivered engine, XeLaTeX with its default Latin Modern text font, reports no missing character.
\begin{prop}\label{Claim:p<n-delta}
	Let $\dF_H= \{A_1, \cdots, A_p\}$ be a greedy partition of an $n$-vertex graph $H$ of size $p$.
	Then, for any $1\leq i<j \leq p,$
	every vertex in $A_j$ has a non-neighbor in $A_i.$ 
    Moreover, we have 
	$p \leq n-\delta(H).$
\end{prop}

\begin{dfn}\label{Dfn:greedy sequence}
An \textit{$(m,p)$-greedy sequence} $A = (a_1, \dots, a_p)$ satisfies that $a_i\in\{1,2,3,4\}$ for $i\in [p],$ $a_i\ge a_{i+1}$ for $i\in [p-1]$, and $\sum_{i=1}^p a_i=m.$ 
Let $a = |\{i\in [p] : a_i = 4\}|$, $b = |\{i\in [p] : a_i \geq 3\}|$, $c = |\{i\in [p] : a_i \geq 2\}|$.
%For an $(m,p)$-greedy sequence $A=\{a_i\}_{i=1}^{p},$
%assuming that the numbers of $4, 3, 2$ and $1's$ in $A$ are $a, b-a, c-b$ and $p-c$ respectively.
Then, the \textit{value} of $A$ is denoted by $f(A):=S_1(A)+S_2(A)+S_3(A),$ where $S_1(A)=b,$ $S_2(A)=mb+a^2-ab-b^2+bc-a-3b,$ and 
\begin{align*}
    S_3(A)=&\sum_{k=3}^{b} a_k \sum_{j=2}^{k-1} (a_j-1) (j-1) +\sum_{k=b+1}^{c} a_k
           \left( \sum_{j=2}^{b}(a_j-1)(j-1) + \sum_{j=b+1}^{k-1}(a_j-1)b \right)+ \\
           &\sum_{k=c+1}^{p} \left( \sum_{j=2}^{b}(a_j-1)(j-1) + \sum_{j=b+1}^{c}(a_j-1)b \right).
\end{align*}
\end{dfn}

\begin{lem}\label{Lem:Bound of 3-clique cover by f(A)}
     If $H$ is a $K_5$-free graph and $\dF_H=\{A_1,\cdots,A_p \}$ is a greedy partition of $H$ of size $p$.  
     Let $A=\{|A_i|\}_{i=1}^{p}$ be the corresponding $(|H|,p)$-greedy sequence of $\dF_H$.
     Then, we have 
	$\coveropt{3}{H}\leq f(A).$
\end{lem}

\begin{proof}
We prove the lemma by constructing a 3-clique cover $\dC$ of $H$ of size at most $f(A).$

Let $a_i = |A_i|$, and let $a,b,c$ be as in Definition~\ref{Dfn:greedy sequence}. 
For $i\in[3],$ let $\dS_i$ be the triangles with three vertices lying in exactly $i$ different $A_j$'s.
First, we can easily see that  $\dC_1:=\bigcup_{i}\{A_i: a_i\geq 3\}$ is a set of cliques that covers all the triangles in $\dS_1.$
Thus, we have $|\dC_1|=|\cup_{i}\{A_i: a_i\geq 3\}|=b=S_1(A).$

For a clique $C$ in $H$ and a vertex set $U\subset V(H)$, recall that $T_U[C]$ is the induced subgraph of $H$ on vertex set $C\cup U',$ where $U'$ is the common neighbor of the vertices in $C$ in $U.$   
Let $\mathcal{I}_1=\{(i,j): 1\le i<j\le p\ \text{and}\ i\le b \},$
    $\mathcal{I}_2=\{(i,j): 1\le i<j\le a\},$ 
    $\mathcal{I}_3=\{(i,j): 1\le i<j\le p,\ i\le b\ \text{and}\ j>a\},$ and 
\begin{align*}
\dC_2=\bigcup\limits_{(i,j)\in\mathcal{I}_1} \{T_{A_i}[v]: v\in A_j\}
      \bigcup\limits_{(i,j)\in\mathcal{I}_2} \{T_{A_j}[v]: v\in A_i\}
      \bigcup\limits_{(i,j)\in\mathcal{I}_3} \{T_{A_i}[u,v]: uv \in E(A_j)\}.
\end{align*}
We claim that $\dC_2$ covers all the triangles in $\dS_2.$ 
Indeed, by %Proposition~\ref{Claim:p<n-delta} and 
the definition of $\dF(H)$, $H[A_{b+1}\cup\cdots\cup A_p]$ is triangle-free.
Thus, the first union set covers all the triangles with one vertex lying in $A_j$ and the other two lying in $A_i$ with $1\leq i<j\leq p$, and the size of the union set is 
\begin{align*}
&\sum_{j=1}^ba_j(j-1)+\sum_{j=b+1}^pa_jb=3\sum_{j=1}^b (j-1)+\sum_{j=1}^a (j-1)+b\sum_{j=b+1}^pa_j \nonumber \\
=&\frac{3}{2}b(b-1)+\frac{1}{2}a(a-1)+b(m-a-3b)=
mb+\frac{a^2}{2}-ab-\frac{3}{2}b^2-\frac {a}{2}-\frac{3}{2}b.
\end{align*} 
Note that the remaining triangles in $\dS_2$ has one vertex lying in $A_i$ and the other two lying in $A_j$ with some $1\le i<j\le p$.
If $|A_i|<|E(A_j)|,$ then we use the second union set to cover them,
and the size of the union set is
$\sum_{i=1}^{a}a_i(a-i)=4\sum_{i=1}^{a}(a-i)=2a^2-2a.$
If $|A_i|\ge |E(A_j)|,$ then we use the third union set to cover them,
and the size of the union set is
\begin{align*}
\sum_{j=a+1}^b\binom{a_j}{2}(j-1)+\sum_{j=b+1}^c\binom{a_j}{2}b
=3\sum_{j=a+1}^b (j-1)+\sum_{j=b+1}^c b=-\frac{3}{2}a^2+ \frac{1}{2}b^2 +bc +\frac 3 2a-\frac 3 2 b.
\end{align*}
Thus, $\dC_2$ is a set of cliques that covers all the triangles in $\dS_2$ (note that $\dC_2$ may contain vertices or edges).
Combining the above three equalities, we get that 
\begin{align*}
    |\dC_2|= mb+a^2-ab-b^2+bc-a-3b=S_2(A).
\end{align*}
    
In a similar way, to cover all the triangles in $\dS_3,$ we define
\begin{align*}
\dC_3=\bigcup\limits_{i,j,k} \{T_{A_i}[u,v]: u\in A_j, v\in A_k, uv\in E(H),\ 1\le i< j<k\le p\}.
\end{align*}
Note that $\dC_3$ may contain edges.
By Proposition~\ref{Claim:p<n-delta}, 
for any $j<k$ and $v\in A_k$, $v$ has at most $a_j-1$ neighbors in $A_j.$   
By the definition of $\dF(H),$ $H[A_{b+1}\cup\cdots\cup A_p]$ is triangle-free and $H[A_{c+1}\cup\cdots\cup A_p]$ is empty, according to the range of $i,j,k$, 
we have
\begin{align}\label{Equ:upper bound on C3}
|\dC_3|\leq &\sum_{k=3}^{b} a_k \sum_{j=2}^{k-1} (a_j-1) (j-1)
+\sum_{k=b+1}^{c} a_k
\left( \sum_{j=2}^{b}(a_j-1)(j-1) + \sum_{j=b+1}^{k-1}(a_j-1)b \right)\nonumber \\
+&\sum_{k=c+1}^{p} \left(  \sum_{j=2}^{b}(a_j-1)(j-1) + \sum_{j=b+1}^{c}(a_j-1)b \right)=S_3(A).
\end{align}

Let $\dC=\dC_1\cup \dC_2 \cup \dC_3\backslash D,$ where $D=\{V(H)\cup E(H)\}.$
Then, by the analysis above and the equations 
$|\dC_1|=S_1(A)$, $|\dC_2|=S_2(A),$ and $|\dC_3|\le S_3(A)$, 
we get that $\dC$ is a $3$-clique cover of $H$ of size at most $f(A)$, which completes the proof of Lemma~\ref{Lem:Bound of 3-clique cover by f(A)}. 
\end{proof}

\end{document}
